\begin{figure}[H]
\centering
\begin{tikzpicture}
    \node[actor, minimum width=5.4cm] (user) at (0,0) {Visitantes y administradores};
    \node[interfaz, minimum width=5.4cm] (presentation) at (0,-2.2) {Presentación pública\\e interfaz CMS};
    \node[servicio, minimum width=5.4cm] (services2) at (0,-4.4) {Servicios CMS: contenido, blog,\\imágenes, respaldos y publicación};
    \node[documento, minimum width=4cm] (contentstore) at (-2.7,-7.1) {Content Store\\JSON, Markdown\\e imágenes};
    \node[datos, minimum width=3cm] (sqlite) at (2.7,-7.1) {SQLite\\operativo};

    \draw[flecha] (user) -- node[etq, right=3pt] {consulta o administra} (presentation);
    \draw[flecha] (presentation) -- node[etq, right=3pt] {solicita operaciones\\(solo el CMS)} (services2);
    \draw[flecha] (services2.south-|contentstore) -- node[etq, left=3pt] {lee y escribe\\contenido} (contentstore);
    \draw[flecha] (services2.south-|sqlite) -- node[etq, right=3pt] {registra\\operación} (sqlite);
    \coordinate (renderroute2) at ([xshift=-1cm]contentstore.west);
    \draw[flecha] (contentstore.west) -- (renderroute2) |- (presentation.west);
    \node[etq, anchor=east] at ([xshift=-0.15cm]$(renderroute2)!0.5!(renderroute2|-presentation.west)$) {contenido\\renderizado};
\end{tikzpicture}
\caption[Capas del CMS en la Iteración 2]{Capas del CMS actualizadas para publicaciones, respaldos y publicación}
\label{fig:cms_iteracion2_extension}
\end{figure}
